\documentclass[11pt,a4paper]{article}

\usepackage[utf8]{inputenc}
\usepackage[T1]{fontenc}
\usepackage[spanish,es-tabla]{babel}
\usepackage{geometry}
\usepackage{enumitem}
\usepackage{listings}
\usepackage[framemethod=tex]{mdframed}
\usepackage{xcolor}
\usepackage{amsmath}
\usepackage{hyperref}

\geometry{margin=2.5cm}

\hypersetup{
  colorlinks=true,
  linkcolor=blue,
  urlcolor=blue
}

\SurroundWithMdFramed[
  nobreak=true,
  linewidth=0.4pt,
  linecolor=black,
  skipabove=6pt,
  skipbelow=6pt,
  innerleftmargin=0pt,
  innerrightmargin=0pt,
  innertopmargin=4pt,
  innerbottommargin=4pt,
]{lstlisting}

\SurroundWithMdFramed[
  nobreak=true,
  linewidth=0.4pt,
  linecolor=black,
  skipabove=4pt,
  skipbelow=4pt,
  innerleftmargin=4pt,
  innerrightmargin=4pt,
  innertopmargin=2pt,
  innerbottommargin=2pt,
]{verbatim}

\lstdefinestyle{pseudocode}{
  basicstyle=\ttfamily\small,
  frame=none,
  breaklines=true,
  columns=fullflexible,
  keepspaces=true,
  showstringspaces=false,
  extendedchars=true,
  literate=
    {á}{{\'a}}1 {é}{{\'e}}1 {í}{{\'i}}1 {ó}{{\'o}}1 {ú}{{\'u}}1
    {Á}{{\'A}}1 {É}{{\'E}}1 {Í}{{\'I}}1 {Ó}{{\'O}}1 {Ú}{{\'U}}1
    {ñ}{{\~n}}1 {Ñ}{{\~N}}1
    {ü}{{\"u}}1 {Ü}{{\"U}}1
    {←}{{$\leftarrow$}}1 {→}{{$\rightarrow$}}1,
}

\lstdefinestyle{ccode}{
  language=C,
  basicstyle=\ttfamily\small,
  frame=none,
  breaklines=true,
  keywordstyle=\color{blue},
  commentstyle=\color{gray},
  stringstyle=\color{teal},
  showstringspaces=false,
  tabsize=4,
}

\title{\textbf{Sistemas Operativos}\\Examen 1}
\author{}
\date{}

\begin{document}
\maketitle

\noindent\textbf{Nombre:}\ \hrulefill\quad
\textbf{Matrícula:}\ \hrulefill\quad
\textbf{Fecha:}\ \hrulefill

\vspace{0.8em}
\noindent Responda con precisión. En diagramas de procesos indique PIDs o roles (padre, hijo, shell, \texttt{nc}, etc.) cuando sea posible. Justifique las respuestas breves.

\vspace{1em}
\hrule
\vspace{1em}

\begin{enumerate}[label=\textbf{\arabic*.}, leftmargin=*]

\item Defina formalmente qué es un \emph{programa} y qué es un \emph{proceso}.

\item Si ejecuta el mismo binario dos veces seguidas, ¿se crea uno o dos procesos? Justifique su respuesta.

\item Explique la diferencia conceptual entre \textbf{concurrencia} y \textbf{paralelismo}. Dé un ejemplo de cada uno en un sistema con un solo núcleo de CPU y otro con varios núcleos.

\item ¿Por qué una configuración de arranque dual (\emph{dual boot}) no se considera una forma de virtualización?

\item ¿Qué tipo de estructura del kernel utiliza el sistema operativo Linux? Mencione una característica distintiva de ese modelo frente a un microkernel.

\item Considere el comando ejecutado en una shell:
\begin{verbatim}
ls | grep ".c"
\end{verbatim}
Responda:
\begin{enumerate}[label=(\alph*), leftmargin=2em]
  \item ¿Cuántos procesos participan (incluya los que crea la shell para el pipeline)?
  \item ¿Qué proceso escribe en la tubería (\emph{pipe})?
  \item ¿Qué proceso lee de la tubería?
  \item ¿Qué tipo de tubería se utiliza: anónima o con nombre?
\end{enumerate}

\item Considere el siguiente programa en C:

\begin{lstlisting}[style=ccode]
#include <stdio.h>
#include <unistd.h>

int main(void)
{
    fork();
    fork();
    execlp("echo", "echo", "hola mundo", (char *)NULL);
    fork();
    return 0;
}
\end{lstlisting}

\begin{enumerate}[label=(\alph*), leftmargin=2em]
  \item Dibuje el árbol de procesos generado desde el \texttt{main} original hasta que todas las ramas terminen o queden bloqueadas en \texttt{exec}/\texttt{fork}.
  \item Indique cuántos procesos existen en el instante en que el último \texttt{fork()} retorna con éxito (antes de que los hijos empiecen a terminar).
  \item ¿Cuántas veces se imprime \texttt{hola mundo} en la salida estándar? Explique qué ocurre con las llamadas a \texttt{execlp} en los distintos procesos.
\end{enumerate}

\item Asuma un servidor basado en el pseudocódigo siguiente (modelo \emph{fork-per-client}, puerto 8080). El proceso padre acepta conexiones; cada cliente lo atiende un hijo. \textbf{No} se muestra \texttt{wait}/\texttt{waitpid} en el pseudocódigo.

\begin{lstlisting}[style=pseudocode]
FUNCIÓN main():
    listen_fd ← crear socket TCP (IPv4)
    asociar listen_fd a "IP local" + puerto 8080
    marcar listen_fd como socket en escucha (cola de conexiones pendientes)

    REPETIR para siempre:
        client_fd ← esperar nueva conexión en listen_fd   // puede bloquear

        pid ← crear proceso hijo (fork)

        SI soy el HIJO (pid == 0):
            cerrar listen_fd en el hijo
            atender_cliente(client_fd)   // al cerrar el cliente, _exit(0)
            // dormir(3 segundos)
            terminar proceso hijo con código 0
        cerrar client_fd en el padre
        // dormir(20 segundos)
    FIN REPETIR
\end{lstlisting}

\noindent En otra terminal, con el servidor ya en ejecución, se lanza:

\begin{verbatim}
for i in 1 2 3; do nc localhost 8080 & done; wait
\end{verbatim}

\noindent El \texttt{wait} final espera a que los tres \texttt{nc} en segundo plano terminen; las tres conexiones se intentan de forma \textbf{concurrente}.

\begin{enumerate}[label=(\alph*), leftmargin=2em]
  \item Tras ejecutar el bucle \texttt{for} anterior (suponga que cada \texttt{nc} se conecta, intercambia datos y cierra la sesión de forma rápida), ¿cuántos procesos del \textbf{programa del servidor} (mismo binario \texttt{miniserv}) existen típicamente en un instante dado mientras los tres clientes aún podrían estar activos? Explique.
  \item Dibuje el \textbf{árbol genealógico} de procesos involucrados: incluya el proceso padre del servidor, sus hijos worker, la shell que ejecutó el \texttt{for} y los procesos \texttt{nc} (puede usar etiquetas \emph{padre servidor}, \emph{hijo 1}, \emph{nc$_1$}, etc.).
  \item En el pseudocódigo, las líneas \texttt{// dormir(3 segundos)} (hijo) y \texttt{// dormir(20 segundos)} (padre) están comentadas. Si \textbf{descomenta ambas}, ¿cómo puede afectar eso la aparición de procesos \textbf{zombie}? Relacione hijo que termina, padre que duerme y ausencia de \texttt{wait}.
  \item Si la línea \texttt{// dormir(20 segundos)} del padre permanece \textbf{comentada}, ¿pueden seguir ejecutándose el servidor y los hijos que atienden clientes? ¿Puede igualmente haber zombies en ese escenario? Argumente.
\end{enumerate}

\item \textbf{Pregunta extra.} Describa qué ocurre durante el \textbf{POST} (\emph{Power-On Self-Test}) en el arranque del sistema. ¿Qué componente lo ejecuta y qué tipo de fallos puede detectar antes de cargar el sistema operativo?

\end{enumerate}

\end{document}
